\documentclass[11pt]{article}

\usepackage{amsmath,amssymb}
\usepackage{xurl}
\usepackage[colorlinks=true,linkcolor=blue,citecolor=blue,urlcolor=blue]{hyperref}

% Shared commands for the notes in docs/paper-gaps/.
% Every note loads this file via % Shared commands for the notes in docs/paper-gaps/.
% Every note loads this file via % Shared commands for the notes in docs/paper-gaps/.
% Every note loads this file via \input{command}. Project-wide notation
% belongs here, never in per-note redefinitions.

% --- Project configuration: edit these three lines for your repository -----
\newcommand{\projectIssueBase}{https://github.com/OWNER/REPO/issues/}
\newcommand{\projectPrBase}{https://github.com/OWNER/REPO/pull/}
\newcommand{\projectDocBase}{https://OWNER.github.io/REPO/docs/}
% ---------------------------------------------------------------------------

\newcommand{\Lean}{\textsc{Lean}}

% A formal declaration name, upright sans, breakable at underscores.
\ExplSyntaxOn
\NewDocumentCommand{\leanid}{m}{
  \ifmmode
    \textsf{\detokenize{#1}}
  \else
    \group_begin:
    \urlstyle{sf}
    \tl_set:Nn \l_tmpa_tl {#1}
    \tl_replace_all:Nnn \l_tmpa_tl {\_} {_}
    \exp_args:NV \path \l_tmpa_tl
    \group_end:
  \fi
}
\ExplSyntaxOff

% Traceability links; use in footnotes, not in the running prose.
\newcommand{\ghissue}[1]{\href{\projectIssueBase#1}{GitHub issue \##1}}
\newcommand{\ghpr}[1]{\href{\projectPrBase#1}{PR \##1}}
\newcommand{\doclink}[1]{\href{\projectDocBase#1.html}{\path{#1}}}

% Domain notation (norms, operators, Dirac kets, ...) for your project's
% mathematics goes below, using \providecommand so a note-local refinement
% stays possible.
\providecommand{\tr}{\operatorname{tr}}

% Verdict marker: \gapnote{<kind>}{<status>}. Kind is the policy
% classification (clarification, local-correction, scope-restriction,
% unfaithful, false-source, open-gap); status is open, wip (elimination
% underway), resolved, or historical. Machine-read by the site index;
% prints a verdict line.
\newcommand{\gapnote}[2]{\par\noindent\textbf{Verdict:} #1 (#2).\par}
. Project-wide notation
% belongs here, never in per-note redefinitions.

% --- Project configuration: edit these three lines for your repository -----
\newcommand{\projectIssueBase}{https://github.com/OWNER/REPO/issues/}
\newcommand{\projectPrBase}{https://github.com/OWNER/REPO/pull/}
\newcommand{\projectDocBase}{https://OWNER.github.io/REPO/docs/}
% ---------------------------------------------------------------------------

\newcommand{\Lean}{\textsc{Lean}}

% A formal declaration name, upright sans, breakable at underscores.
\ExplSyntaxOn
\NewDocumentCommand{\leanid}{m}{
  \ifmmode
    \textsf{\detokenize{#1}}
  \else
    \group_begin:
    \urlstyle{sf}
    \tl_set:Nn \l_tmpa_tl {#1}
    \tl_replace_all:Nnn \l_tmpa_tl {\_} {_}
    \exp_args:NV \path \l_tmpa_tl
    \group_end:
  \fi
}
\ExplSyntaxOff

% Traceability links; use in footnotes, not in the running prose.
\newcommand{\ghissue}[1]{\href{\projectIssueBase#1}{GitHub issue \##1}}
\newcommand{\ghpr}[1]{\href{\projectPrBase#1}{PR \##1}}
\newcommand{\doclink}[1]{\href{\projectDocBase#1.html}{\path{#1}}}

% Domain notation (norms, operators, Dirac kets, ...) for your project's
% mathematics goes below, using \providecommand so a note-local refinement
% stays possible.
\providecommand{\tr}{\operatorname{tr}}

% Verdict marker: \gapnote{<kind>}{<status>}. Kind is the policy
% classification (clarification, local-correction, scope-restriction,
% unfaithful, false-source, open-gap); status is open, wip (elimination
% underway), resolved, or historical. Machine-read by the site index;
% prints a verdict line.
\newcommand{\gapnote}[2]{\par\noindent\textbf{Verdict:} #1 (#2).\par}
. Project-wide notation
% belongs here, never in per-note redefinitions.

% --- Project configuration: edit these three lines for your repository -----
\newcommand{\projectIssueBase}{https://github.com/OWNER/REPO/issues/}
\newcommand{\projectPrBase}{https://github.com/OWNER/REPO/pull/}
\newcommand{\projectDocBase}{https://OWNER.github.io/REPO/docs/}
% ---------------------------------------------------------------------------

\newcommand{\Lean}{\textsc{Lean}}

% A formal declaration name, upright sans, breakable at underscores.
\ExplSyntaxOn
\NewDocumentCommand{\leanid}{m}{
  \ifmmode
    \textsf{\detokenize{#1}}
  \else
    \group_begin:
    \urlstyle{sf}
    \tl_set:Nn \l_tmpa_tl {#1}
    \tl_replace_all:Nnn \l_tmpa_tl {\_} {_}
    \exp_args:NV \path \l_tmpa_tl
    \group_end:
  \fi
}
\ExplSyntaxOff

% Traceability links; use in footnotes, not in the running prose.
\newcommand{\ghissue}[1]{\href{\projectIssueBase#1}{GitHub issue \##1}}
\newcommand{\ghpr}[1]{\href{\projectPrBase#1}{PR \##1}}
\newcommand{\doclink}[1]{\href{\projectDocBase#1.html}{\path{#1}}}

% Domain notation (norms, operators, Dirac kets, ...) for your project's
% mathematics goes below, using \providecommand so a note-local refinement
% stays possible.
\providecommand{\tr}{\operatorname{tr}}

% Verdict marker: \gapnote{<kind>}{<status>}. Kind is the policy
% classification (clarification, local-correction, scope-restriction,
% unfaithful, false-source, open-gap); status is open, wip (elimination
% underway), resolved, or historical. Machine-read by the site index;
% prints a verdict line.
\newcommand{\gapnote}[2]{\par\noindent\textbf{Verdict:} #1 (#2).\par}


\title{Model Note: Comparing a Cited Argument with a Formal Statement}
\author{}
\date{\today}

\begin{document}
\maketitle
\gapnote{unfaithful}{open}

% This file is a model for notes in docs/paper-gaps/.
% When making a real note, replace the abstract toy assertion below by the actual
% mathematical assertion, source citation, issue link, and formal declarations.
% The visible document should remain a coherent mathematical note, not a form.
% If the issue is a scalar correction, a missing term, a counterexample, or a
% different formal-library route, specialize the middle sections accordingly.
% If a Mathlib theorem or construction replaces a step from the paper, explain
% that theorem in ordinary mathematics before naming the formal declaration.
% In a real note, the first paragraph should identify the cited assertion and
% the part of the argument in which it is used.  Put issue links, pull requests,
% source files, line ranges, labels, and formal declaration names in footnotes
% when they are only traceability data.

\section{The assertion}

This note concerns a hypothetical proposition, of the sort one might encounter
in a cited argument such as \cite{ModelSource}.  The proposition is used
later in the same argument.

Let \(\mathcal H\) denote the hypotheses available at a certain point of the
argument, and let \(C\) denote the conclusion needed there.  Suppose the cited
source states the implication
\begin{align}
  \mathcal H \Longrightarrow C.
  \tag{\path{thm:model-source}}
\end{align}
The notation in the displayed implication has now been fixed: \(\mathcal H\) is
the available hypothesis set, and \(C\) is the conclusion required by the later
argument.

\section{The point at issue}

Assume that the proof of the displayed implication invokes a lemma which, in the
special case being used, proves
\begin{align}
  \mathcal H \wedge \mathcal K \Longrightarrow C,
  \tag{\path{lem:model-input}}
\end{align}
where \(\mathcal K\) is an additional hypothesis.  If the surrounding argument
does not prove \(\mathcal K\), then the displayed implication
\(\mathcal H\Rightarrow C\) has not been justified by that proof.  The
mathematical question is whether \(\mathcal K\) follows from \(\mathcal H\), must
be added to the statement, or can be avoided by a different argument.

Thus the proof supplies the desired conclusion only after the additional
hypothesis \(\mathcal K\) has been accounted for.

\section{The corrected statement}

In this model, suppose that \(\mathcal K\) is not known to follow from
\(\mathcal H\), but that the later proof only uses the theorem in situations
where \(\mathcal K\) is available.  The corrected statement is then
\begin{align}
  \mathcal H \wedge \mathcal K \Longrightarrow C.
\end{align}
This is a strengthening of the hypotheses, not a derivation of \(\mathcal K\)
from the original assumptions.

If the formal development proves this corrected statement, the note should state
that fact in prose and identify the exact formal declaration in a footnote.  For
example, the declaration might be recorded as \leanid{modelCorrected} in
\doclink{Project/Path/To/File}.  The formal declaration is evidence for the
corrected mathematical statement; it is not a substitute for explaining the
statement.

The same principle applies if the formal proof uses a theorem from a library
which is not present in the cited source.  The note should first state that
theorem as a mathematical assertion, then explain how its hypotheses are verified
in the present notation, and only then identify the corresponding formal
declaration.

\section{Conclusion}

The verdict in this model is that the original proof uses an additional
hypothesis \(\mathcal K\).  The corrected statement is sufficient for the
surrounding argument only in the cases where \(\mathcal K\) is supplied.  Without
such a derivation or a different proof of \(C\), the original implication
\(\mathcal H\Rightarrow C\) remains unjustified.

\bibliographystyle{alpha}
\bibliography{references}

\end{document}
